This work presented EPICA, an integrated platform for Langmuir probe diagnostics that combines probe bias generation, signal conditioning, data acquisition, and embedded processing within a modular architecture. The system was specifically developed to address the scalability challenges associated with large Langmuir probe arrays, such as those foreseen for MITICA.

Experimental validation confirmed compliance with the measurement requirements defined for MITICA operation. The proposed architecture successfully demonstrated high-resolution voltage and current acquisition and the capability to perform complete Langmuir probe measurements.

Unlike conventional implementations based on centralized laboratory instrumentation, EPICA integrates the complete measurement chain within each probe module, significantly reducing wiring complexity, footprint, and system-level integration effort. The combination of distributed bias generation, FPGA-based processing, and compatibility with ITER CODAC provides a scalable solution for large plasma diagnostic systems.

The modular organization of the platform enables straightforward scaling towards the 212 measurement points foreseen for MITICA while preserving synchronization, centralized supervision, and measurement performance.

Future developments will focus on the deployment of the complete MARTe2 software framework and on extended validation of the proposed architecture in ongoing plasma experiments. These activities will provide further assessment of the system performance under realistic operating conditions and support future large-scale deployments.